\chapter{Stack tecnologico}
\label{chap:stack}

Questa sezione motiva le scelte tecnologiche adottate per ciascun componente, descrivendo
le alternative considerate e i criteri di selezione.

\section{Backend: FastAPI}
\label{sec:stack_backend}

Il backend \`e sviluppato in \textbf{Python 3.9} con il framework \textbf{FastAPI}. Le
alternative principali considerate erano Django e Flask.

\textbf{Django} offre una soluzione completa e matura (ORM, admin panel, auth integrata),
ma la sua architettura monolitica mal si adatta a un sistema distribuito in cui la logica
applicativa \`e intenzionalmente distribuita tra microservizi. La presenza di molte
funzionalit\`a non necessarie aumenta anche la superficie di attacco.

\textbf{Flask} \`e estremamente leggero e flessibile, ma richiede l'integrazione manuale
di validazione dei dati, documentazione API e gestione asincrona, aggiungendo complessit\`a
senza benefici tangibili rispetto a FastAPI.

\textbf{FastAPI} \`e stato scelto per:
\begin{itemize}[leftmargin=*]
    \item Supporto nativo alla programmazione asincrona (\texttt{async/await}),
          fondamentale per I/O concorrente (database, S3, chiamate esterne).
    \item Validazione automatica degli input tramite \textbf{Pydantic v2}: ogni
          schema di request/response \`e definito come modello Python con type hints.
    \item Generazione automatica della documentazione OpenAPI/Swagger UI
          all'endpoint \texttt{/docs}.
    \item Compatibilit\`a nativa con \textbf{SQLAlchemy 2} e il pattern
          dependency injection tramite \texttt{Depends}.
    \item Gestione centralizzata degli errori tramite \texttt{exception\_handler}:
          gli errori di validazione producono risposte normalizzate.
\end{itemize}

Il server WSGI/ASGI utilizzato \`e \textbf{Uvicorn} con driver \textbf{psycopg3}
(\texttt{psycopg[binary]}) per la connessione asincrona a PostgreSQL.

\section{Database relazionale: PostgreSQL 16}
\label{sec:stack_db}

Per la gestione dei dati strutturati \`e stato adottato \textbf{PostgreSQL 16}.
L'alternativa principale era MySQL; SQLite \`e stato scartato a priori per assenza
di supporto alla concorrenza ad alto carico.

\textbf{MySQL} \`e una soluzione diffusa e stabile, con buone prestazioni nelle operazioni
di lettura, ma con funzionalit\`a pi\`u limitate nella gestione di tipi complessi e nella
ricerca full-text nativa.

La scelta di \textbf{PostgreSQL} \`e motivata da:
\begin{itemize}[leftmargin=*]
    \item Supporto nativo al tipo \textbf{JSONB} per la colonna \texttt{details} della
          tabella \texttt{agent\_operations}, che accoglie strutture arbitrarie per ogni
          tipo di evento agentico.
    \item Generazione di UUID tramite l'estensione \textbf{pgcrypto}
          (\texttt{gen\_random\_uuid()}), evitando la dipendenza da sequenze incrementali.
    \item Supporto a \textbf{transazioni ACID} complete, essenziale per garantire la
          consistenza delle operazioni di filing che toccano pi\`u tabelle.
    \item Ricca espressivit\`a SQL: \texttt{CHECK} constraint, \texttt{ON DELETE CASCADE},
          indici compositi e copertura.
    \item \textbf{SQLAlchemy 2} fornisce un ORM maturo e type-safe con lazy loading
          controllato e session management esplicito.
\end{itemize}

\section{Object storage: SeaweedFS}
\label{sec:stack_s3}

Per la gestione dei file binari (documenti e \texttt{.eml}) \`e stato adottato
\textbf{SeaweedFS}, un object store S3-compatibile open source.

Le alternative considerate erano:
\begin{description}[leftmargin=*]
    \item[S3 gestito (AWS, GCP, Azure)] Lo standard industriale in termini di affidabilit\`a
          e scalabilit\`a, ma introduce dipendenza da provider esterno e costi variabili
          difficilmente predicibili in fase di sviluppo.
    \item[MinIO] Semplice da integrare e S3-compatibile, ma il progetto ha modificato la
          propria licenza (AGPL) rendendola incompatibile con molti scenari SaaS.
    \item[Ceph] Soluzione enterprise con alta resilienza e funzionalit\`a avanzate, ma con
          complessit\`a operativa elevata (richiede cluster dedicato, expertise specializzata).
    \item[Garage] Alternativa recente orientata alla resilienza distribuita e geografica,
          ma con ecosistema ancora immaturo.
\end{description}

\textbf{SeaweedFS} \`e stato scelto per il bilanciamento tra semplicità operativa, prestazioni
e compatibilit\`a S3. L'interfaccia standard S3 (tramite boto3) garantisce la possibilit\`a di
migrare verso un provider gestito senza modifiche all'applicazione.

Un vantaggio architetturale chiave \`e l'utilizzo di \textbf{presigned URL}: il dato binario
non transita mai attraverso il backend, che agisce solo come piano di controllo. Questo riduce
il traffico verso il backend, abbassa la latenza percepita dall'utente e disaccoppia lo scaling
del backend da quello dello storage.

\section{Estrazione testo: Apache Tika}
\label{sec:stack_tika}

Per l'estrazione di testo da documenti binari eterogenei \`e stato adottato
\textbf{Apache Tika}, eseguito come container REST.

L'alternativa sarebbe stata l'integrazione di librerie specifiche per formato:
\texttt{pypdf2}/\texttt{pdfplumber} per PDF, \texttt{python-docx} per DOCX,
\texttt{openpyxl} per XLSX. Questo approccio richiederebbe la manutenzione di
parser multipli con comportamenti disomogenei, soprattutto per edge case (PDF con
encoding non standard, documenti protetti, formati ibridi).

\textbf{Apache Tika} offre un'interfaccia unificata per oltre 1000 tipi di file,
semplificando la pipeline di elaborazione. Nel sistema viene invocato tramite HTTP PUT
con timeout di 30~s per prevenire hang su file corrotti o documenti molto grandi. I formati
supportati rilevanti per il dominio sono: PDF, DOCX, XLSX, ODT, TXT, RTF.

\section{Ricerca: OpenSearch}
\label{sec:stack_opensearch}

Per le funzionalit\`a di ricerca avanzata \`e stato adottato \textbf{OpenSearch 2.18.0},
motore di ricerca open source derivato da Elasticsearch.

La ricerca full-text nativa di PostgreSQL (\texttt{tsvector}/\texttt{tsquery}) \`e stata
scartata per la difficolt\`a di combinare ricerca testuale e ricerca vettoriale in modo
efficiente su dataset di grandi dimensioni. Elasticsearch sarebbe stato equivalente ma
ha modificato la propria licenza verso un modello misto non completamente open source.

\textbf{OpenSearch} offre:
\begin{itemize}[leftmargin=*]
    \item \textbf{Full-text search} con analizzatore \texttt{standard} sul campo \texttt{text}.
    \item \textbf{Ricerca vettoriale k-NN} tramite \textsc{HNSW} (Hierarchical Navigable
          Small World) con metrica \texttt{cosinesimil} su embedding a 384 dimensioni.
    \item \textbf{Indicizzazione per tenant}: ogni utente ha un proprio indice
          \texttt{documents\_\{user\_id\_no\_hyphens\}} che garantisce isolamento e
          possibilit\`a di scaling indipendente.
\end{itemize}

Gli embedding sono generati dal modello
\textbf{paraphrase-multilingual-MiniLM-L12-v2} di \texttt{sentence-transformers},
scelto per il supporto multilingue (rilevante in contesti aziendali italiani ed europei)
e per le dimensioni contenute del vettore (384 vs.\ 768 o 1536 di modelli pi\`u grandi).

\section{Classificazione: Atomic Agents e Instructor}
\label{sec:stack_agent}

L'agente di classificazione \`e implementato con il framework \textbf{Atomic Agents},
che fornisce un'architettura modulare per agenti LLM con schema di input/output
tipizzato tramite Pydantic.

La libreria \textbf{instructor} \`e usata come adattatore per forzare l'output del modello
LLM a conformarsi allo schema Pydantic \texttt{FilingOutput}. Questo approccio elimina
la necessit\`a di parsing manuale dell'output e garantisce che risposte strutturalmente
non conformi vengano automaticamente rigenerate.

Il client LLM \`e configurato tramite \textbf{API OpenAI-compatible} puntando a un
endpoint configurabile tramite variabile d'ambiente \texttt{BASE\_URL}, rendendo il
sistema indipendente dal provider. Il modello di default \`e
\textbf{claude-sonnet-4-6} (Anthropic), accessibile tramite endpoint compatibile.

\section{Email polling: IMAPClient e APScheduler}
\label{sec:stack_email}

La componente di polling email usa:
\begin{itemize}[leftmargin=*]
    \item \textbf{IMAPClient}: libreria Python per il protocollo IMAP (RFC 3501) con
          supporto nativo a SSL/TLS e IMAP UID (meccanismo di watermarking).
    \item \textbf{APScheduler 3.11.2}: scheduler in-process per l'esecuzione periodica
          del loop di polling, con supporto per intervalli configurabili per account.
\end{itemize}

\section{Frontend: React e Vite}
\label{sec:stack_frontend}

L'interfaccia utente \`e una \textbf{Single Page Application} sviluppata con \textbf{React 18}
e il build tool \textbf{Vite 5}, scelto per i tempi di build significativamente inferiori
rispetto a \texttt{create-react-app}/Webpack.

Le scelte di librerie sono volutamente minimaliste: \texttt{react-router-dom} per il routing,
\texttt{lucide-react} per le icone, CSS vanilla per lo stile. L'assenza di state manager
esterno (Redux, Zustand) \`e una scelta deliberata: la complessit\`a dello stato dell'applicazione
non giustifica il costo di onboarding di tali librerie.

La build \`e prodotta con un Dockerfile \textbf{multi-stage}: uno stage Node per la compilazione
e uno stage Nginx per servire i file statici risultanti, riducendo la dimensione dell'immagine
finale e la superficie di attacco.
